\documentclass[10pt,twocolumn]{article}
\usepackage[T1]{fontenc}
\usepackage[utf8]{inputenc}
\usepackage{lmodern}
\usepackage[margin=1.8cm,top=2.4cm,bottom=2cm,columnsep=0.6cm]{geometry}
\usepackage{fancyhdr}
\usepackage{titlesec}
\usepackage{booktabs}
\usepackage{amsmath,amssymb,amsfonts}
\usepackage{microtype}
\usepackage{xcolor}
\usepackage{mdframed}
\usepackage{parskip}
\usepackage{enumitem}
\usepackage{hyperref}

\definecolor{rulered}{RGB}{180,30,30}
\definecolor{boxgray}{RGB}{245,245,245}
\definecolor{keywordgray}{RGB}{100,100,100}

\pagestyle{fancy}
\fancyhf{}
\fancyhead[L]{\textsc{\small Claude Mag}}
\fancyhead[C]{\small\textit{Machine Learning Research Digest}}
\fancyhead[R]{\small 2026-11-09}
\fancyfoot[C]{\small\thepage}
\renewcommand{\headrulewidth}{0.8pt}

\titleformat{\section}{\normalsize\bfseries\color{rulered}}{}{0pt}{}%
  [\vspace{-4pt}\color{rulered}\rule{\columnwidth}{0.4pt}]
\titlespacing{\section}{0pt}{8pt}{4pt}

\hypersetup{colorlinks=true,urlcolor=rulered,linkcolor=black}

\begin{document}
\setlength{\parindent}{0pt}
\setlength{\parskip}{4pt}

{\small\color{keywordgray}\textbf{Keywords:}
  medical AI \textbullet{}
  graph memory \textbullet{}
  clinical dialogue \textbullet{}
  patient state tracking}\par\vspace{4pt}

{\LARGE\bfseries\raggedright
  MedGraph: Maintaining Patient State with Graph Memory
  for Long-Term Medical Dialogue}\par\vspace{4pt}

{\large\itshape\raggedright
  Three-Layer Hierarchical Graph Memory Eliminates Outdated State,
  Cross-Patient Contamination, and Missing Safety Facts, Gaining
  10.5 Points on MedMemoryBench}\par\vspace{6pt}

{\small\textbf{By} Donghua Zhang, Zizheng Wang, Ze Dong, Yiding Feng,
  Siyu Qian}\par\vspace{2pt}

{\small\color{keywordgray}arXiv:2611.08433 \textbullet{}
  cs.CL / cs.AI \textbullet{} submitted 6 Nov 2026}
\par\vspace{2pt}\noindent{\color{rulered}\rule{\columnwidth}{0.6pt}}\vspace{6pt}

Medical dialogue assistants fail in longitudinal settings because
flat memory systems accumulate outdated states, mix facts across
patients, and silently drop safety information. MedGraph introduces
a three-layer hierarchical graph memory that makes clinical state
transitions explicit and unconditionally prepends safety facts to
every query. On MedMemoryBench it exceeds the strongest of five
baselines by 10.4--10.5 percentage points with both GPT and Qwen
backbones.

\begin{mdframed}[backgroundcolor=boxgray,linecolor=rulered,linewidth=1pt,
    innerleftmargin=6pt,innerrightmargin=6pt,
    innertopmargin=6pt,innerbottommargin=6pt]
\textbf{\small AT A GLANCE}\par\vspace{4pt}
\begin{description}[leftmargin=0pt,labelsep=4pt,itemsep=2pt,parsep=0pt,
    font=\small\bfseries]
  \item[Problem:] Flat memory systems for clinical dialogue produce
    outdated state, cross-patient contamination, and silently missed
    safety facts across multi-visit conversations.
  \item[Why it matters:] LLMs used as clinical assistants in
    longitudinal care must maintain accurate, up-to-date patient
    state; errors can cause harmful medication or diagnosis decisions.
  \item[Method:] Three-layer memory: raw transcripts (L1), a typed
    clinical fact graph with temporal validity tags (L2), and a
    patient profile unconditionally attached to every query (L3).
  \item[Result:] 10.4--10.5 pp above the strongest of five baselines
    on MedMemoryBench with GPT and Qwen backbones.
  \item[Evidence:] MedMemoryBench (multi-visit clinical dialogues);
    baselines: flat vector, sliding window, RAG, Zep, MemoryOS.
\end{description}
\end{mdframed}

\section{The Big Picture}

Large language models deployed as clinical dialogue assistants need
to maintain accurate knowledge of a patient's medical state across
many visits over weeks or months. Existing memory approaches --- flat
vector stores, sliding-window context, retrieval-augmented generation
--- treat memory as a bag of text fragments. This fails in three
recurring ways identified by the authors through systematic error
analysis of five baseline systems on longitudinal medical dialogues:

\begin{enumerate}[itemsep=2pt,parsep=0pt]
  \item \textbf{Outdated state:} A resolved symptom is reported as
    currently active because an older transcript fragment scores
    higher in retrieval than the later discharge note.
  \item \textbf{Cross-patient contamination:} Facts from a different
    patient's earlier session pollute the current patient's context
    when embedding similarity overrides patient identity.
  \item \textbf{Missing safety facts:} A known allergy or
    contraindication is not retrieved because it was not in the
    top-$k$ most similar fragments for the current query.
\end{enumerate}

MedGraph introduces structure that makes all three failures impossible
by design: typed temporal graph updates prevent outdated state; strict
per-patient graph isolation prevents contamination; and mandatory
safety-block attachment prevents silent omission.

\section{Why Existing Approaches Fall Short}

\textbf{Flat vector memory} embeds all clinical utterances and
retrieves the top-$k$ most similar. It does not distinguish whether
a fact is currently active or resolved, and similarity can produce
facts from other patients if patients share medical vocabulary.

\textbf{Sliding-window context} keeps the most recent $n$ tokens.
Long-term facts (a years-old allergy, a chronic condition diagnosed
at first visit) fall out of the window and are lost.

\textbf{Retrieval-augmented generation (RAG)} retrieves relevant
documents at query time. It improves recall of historical facts but
cannot distinguish the current clinical state from historical states
when multiple conflicting observations exist in the corpus.

\textbf{Zep (entity memory)} maintains entity states but uses
flat entity records without typed temporal transitions; it conflates
active and resolved states.

\textbf{MemoryOS} provides structured memory with hierarchical
layers but does not specialize for clinical entity types or enforce
mandatory safety-fact prepending.

\section{Core Mechanism: Three-Layer Memory}

\textbf{Layer 1 --- Raw dialogues.}
Full session transcripts are stored verbatim, enabling exact
retrieval of clinical reasoning chains and direct-quotation evidence
when needed. This layer is searched by traditional BM25 retrieval
for specific historical passages.

\textbf{Layer 2 --- Clinical fact graph.}
A typed knowledge graph over entities (symptoms, diagnoses,
medications, lab results, procedures) with typed relations
(has-symptom, prescribed, diagnosed-with, resolved, contraindicated-with).
Each entity node carries a temporal validity tag: \textit{active},
\textit{resolved}, or \textit{historical}. Graph construction and
update are performed by a clinical NLP extraction pipeline applied to
each new session.

\textbf{Layer 3 --- Patient profile (safety block).}
A compact structured record per patient containing: recorded allergies,
contraindications, chronic conditions, emergency contacts, and
medication alerts. This block is \emph{unconditionally prepended} to
every query context, bypassing the retrieval mechanism entirely.
No fact in this block can be missed because retrieval does not decide
whether it is included.

\textbf{Graph isolation.} Each patient has a fully separate graph
instance and a separate patient profile. Cross-patient contamination
is architecturally impossible.

\section{Technical Formulation}

Let $\mathcal{G}_p = (V_p, E_p)$ be the clinical fact graph for
patient $p$, where each node $v \in V_p$ is a clinical entity with
attributes $(\text{type}, \text{value}, \text{validity})$ and each
edge $e = (v_i, v_j, r)$ is a typed relation.

When a new session transcript $D_t$ arrives for patient $p$, the
extraction pipeline produces a set of entity mentions
$\hat{V}_t$ and relation mentions $\hat{E}_t$. The update procedure
applies type-specific rules:
\begin{align*}
  &\text{if } v \in \hat{V}_t \text{ and } v \notin V_p:
    \quad \text{add } v \text{ with validity} = \textit{active},\\
  &\text{if } v \in \hat{V}_t \text{ and } v \in V_p
    \text{ and ``resolved'' signal:}
    \quad \text{set validity} \to \textit{resolved},\\
  &\text{if } v \in \hat{V}_t \text{ and } v \in V_p
    \text{ and no state change:}
    \quad \text{update timestamp only}.
\end{align*}

The ``resolved'' signal is detected by a small classifier trained on
clinical resolution phrases (``symptom resolved'', ``discontinued'',
``normal finding'', etc.).

At query time, the context assembled for the LLM backbone is:
\[
  \mathcal{C} = \underbrace{\text{Profile}_p}_{\text{L3, mandatory}}
    \oplus \underbrace{\text{Graph}_p^{\text{active}}}_{\text{L2, active facts}}
    \oplus \underbrace{\text{RAG}(\mathcal{D}, q)}_{\text{L1, retrieved}},
\]
where $\text{Graph}_p^{\text{active}}$ is the subgraph of currently
active entities serialized as structured text, and
$\text{RAG}(\mathcal{D}, q)$ is BM25-retrieved transcript fragments
relevant to query $q$.

\section{Learning or Inference Procedure}

MedGraph uses no end-to-end learning. Its components are:
\begin{itemize}[itemsep=2pt,parsep=0pt]
  \item A clinical NER/RE pipeline (pre-trained, fine-tuned on MIMIC-IV
    notes) for entity and relation extraction.
  \item A resolution-signal classifier (a small BERT-based model
    fine-tuned on clinical resolution annotation).
  \item An arbitrary LLM backbone (tested with GPT-4o-mini and
    Qwen2.5-72B) for the final answer generation.
\end{itemize}
The NLP pipeline and classifier are frozen at inference time.

\section{What the Guarantee Says}

No formal guarantee is claimed. The empirical claim is superior
clinical accuracy across all five baselines on MedMemoryBench with
both LLM backbones. The error analysis identifies three failure modes
and shows that MedGraph's design eliminates each by construction;
the reduction in those specific error types is verified per-category.

\section{Experimental Findings}

\begin{itemize}[itemsep=2pt,parsep=0pt]
  \item \textbf{Overall accuracy (GPT-4o-mini):} MedGraph
    $+10.5$ pp above best baseline (Zep).
  \item \textbf{Overall accuracy (Qwen2.5-72B):} MedGraph
    $+10.4$ pp above best baseline.
  \item \textbf{Outdated-state errors:} Reduced by 68\% relative
    to Zep, confirming typed temporal update effectiveness.
  \item \textbf{Safety-fact recall:} 100\% safety-fact recall
    (by mandatory prepending), vs.\ 73--82\% for retrieval-based
    baselines.
  \item \textbf{Cross-patient contamination rate:} 0\% (by graph
    isolation), vs.\ 3--7\% for vector-memory baselines.
\end{itemize}

\section{Ablations and Interpretation}

\textbf{Without mandatory safety block:} Dropping the unconditional
prepend and relying on retrieval for safety facts reduces safety-fact
recall from 100\% to 79\%, confirming that retrieval alone is
insufficient.

\textbf{Without typed temporal updates:} Using simple entity overwrite
(new mention always replaces old) increases outdated-state errors by
3.2$\times$, as resolved symptoms return when re-mentioned in a summary
context.

\textbf{Without graph isolation:} Merging all patients into a single
graph increases cross-patient contamination to 4.1\%, confirming that
per-patient isolation is necessary.

\textbf{L1 retrieval necessity:} Removing L1 (raw transcript retrieval)
reduces overall accuracy by 4.3 pp, showing that verbatim historical
passages provide useful evidence beyond what the structured graph captures.

\section*{Reference}

Donghua Zhang, Zizheng Wang, Ze Dong, Yiding Feng, Siyu Qian.
\textbf{MedGraph: Maintaining Patient State with Graph Memory for
Long-Term Medical Dialogue.}
arXiv:2611.08433.
\url{https://arxiv.org/abs/2611.08433}

\end{document}
